\begin{table}[tbp]
\centering
\caption{\textbf{Six ways to learn or bound the follow-up detection
multiplier.} What each design supplies and the assumption on which it turns.
$M = d_1/d_0$ is defined for a stated condition, finding class or procedure;
the designs are not ranked, because they answer different questions with
different information requirements. Details and the remaining assumptions
for each are in Section~\ref{sec:operational}; the first two identifying
results are Propositions~\ref{prop:ref} and \ref{prop:twoproc} in
Appendix~\ref{app:id}.}
\label{tab:designs}
\small
\begin{tabular}{p{0.30\textwidth}p{0.30\textwidth}p{0.32\textwidth}}
\toprule
Design & What it supplies & Critical assumption \\
\midrule
Reference observation under a common measurement rule & $\Rlat$, hence $M = \Robs/\Rlat$ & Common sensitivity in the two history groups; no false positives \\
High-sensitivity reference procedure & Anchor for $r$, hence $d = q/r$ & Known or credibly near-perfect sensitivity for the target condition \\
Independent reference adjudication & $d_0$, $d_1$ directly & Reference status credible enough to stand in for the condition \\
Two distinct procedures & $r$, $d$ within a history group & Conditionally independent failures; stable sensitivity; no false positives \\
Randomized or quasi-random attention shift & Relative detection; tighter bounds & Exclusion restriction: the shift does not alter or select on the condition \\
Procedure-based bound & $\Mbar$, hence $\Rlat \ge \Robs/\Mbar$ & A defensible model linking procedures to detection probability \\
\bottomrule
\end{tabular}
\end{table}
